% ============================================================================
% PeerJ Computer Science submission draft — structure-compliant equivalent.
% NOTE: the official peerj.cls / peerj.bst ship with the PeerJ Overleaf
% template (https://www.overleaf.com/latex/templates). At submission time,
% import this body into the official template and swap the preamble
% (TODO_FOR_AUTHOR #1). Section names, author-year citations, line numbers,
% and the declaration block below follow the official structure.
% ============================================================================
\documentclass[10pt]{article}
\usepackage[margin=2.5cm]{geometry}
\usepackage[mathlines]{lineno}
\usepackage{amsmath,amssymb}
\usepackage{graphicx}
\usepackage{booktabs}
\usepackage{multirow}
\usepackage{array}
\newcolumntype{L}[1]{>{\raggedright\arraybackslash}p{#1}}
\usepackage{xcolor}
\usepackage{algorithm}
\usepackage{algpseudocode}
\usepackage{tikz}
\usetikzlibrary{arrows.meta,positioning,fit,backgrounds,calc}
\usepackage{caption}
\captionsetup{font=small,labelfont=bf}
\usepackage[round,authoryear,sort&compress]{natbib}
\setlength{\bibsep}{2pt}
\usepackage[colorlinks=true,linkcolor=blue!60!black,citecolor=blue!60!black,urlcolor=blue!60!black]{hyperref}
\sloppy
\setlength{\emergencystretch}{3em}
\linenumbers

\title{\bfseries From Embodied Experience to Persistent Experience-Derived State: An Auditable Unified-Graph Architecture and a Preregistered Boundary Analysis of Activation-Based Access}
% [TODO-FOR-AUTHOR: affiliations, ORCID, corresponding author — must match
%  the online submission metadata EXACTLY (PeerJ requirement).]
\author{Kailin Yan\thanks{Corresponding author: \texttt{2024015943@st.cupk.edu.cn} [TODO-FOR-AUTHOR: confirm affiliation wording and ORCID]}}
\date{}

\begin{document}
\maketitle

\begin{abstract}
Embodied agents accumulate experience far faster than they can write it back as persistent state: action histories remain logs. We present \textsc{Fascinator} (FAS), a cognitive-architecture prototype that keeps semantic, episodic, and procedural content on a single typed, weighted graph and writes experience back as relational structure through one evidence-weighted attribution pipeline. The paper contributes an auditable memory architecture together with a preregistered evaluation protocol that delimits when spreading activation---the architecture's access mechanism---adds value over flat retrieval. In an information-asymmetric beacon task, persistent graph context yields 20/20 task completions versus 0/20 without context (McNemar $p{=}2\times10^{-6}$; the 0/20 is a closed-book design's built-in existence verification); an activation-free embedding-retrieval control reproduces the full benefit (20/20). Persistent experience-derived state is therefore supported as behaviorally effective, while the necessity of activation is not. A context-removal control experiment establishes that persistence itself is causally operative: after graph reset, goal resumption drops from 10/10 to 0/10 (exact McNemar $p{=}0.00195$, the minimum attainable at $n{=}10$) and repeated failures on an invalidated goal rise from a mean of 2.8 to 7.7 (one-sided Wilcoxon $p{=}0.00098$); whenever the experience stream remains re-readable, history and retrieval baselines match FAS---so the causal role attaches to persistence itself, while the independent contribution of graph structure and activation remains open. A preregistered mechanism campaign ($4\times50$ paired seeds) rejects activation's incremental value over a closed-form flat baseline ($\Delta$JCG${<}0$, Holm $p{=}1.5\times10^{-9}$, 0/200 reversals), isolating one exception: multi-input convergence through a shared hub (FAS 30/30 \emph{vs} flat 0/30). An 84-cell parameter sweep maps this exception---FAS strictly beats flat in 11 cells (weak edge weights, shallow paths), loses in 9 (strong weights, deeper paths), and ties in 64. A second, baseline-controlled campaign (508 archived runs; 390 repair runs archived separately) measures the trade-off directly: under neutral instructions in the short-history regime, the persistent state resumes an abandoned goal in 9/10 runs where full-history LLM baselines do so in 0/10 (longer, failure-containing streams: baselines catch up at 10/10), and the same persistence repeats an invalidated goal 10.5 times on average where the baseline pivots after 1.0. All negative results, two assembly audits, and the full protocol are reported.
\end{abstract}

\section{Introduction}
\label{sec:intro}

As soon as an embodied agent acts for more than a few minutes, it produces a stream of experience: actions, observations, successes, failures. What happens to that stream separates two classes of system. In one class, experience accumulates as logs---available for reflection prompts or retrieval augmentation, but never re-organized into structure that changes future behavior. In the other, experience is consolidated: concrete episodes give way to general relational knowledge that persists, is re-activated in relevant situations, and measurably participates in later decisions. Whether or not the system's core includes a large language model (LLM), the gap between logs and consolidation is a central engineering problem for long-lived agents.

Existing positions are familiar. Cognitive architectures such as ACT-R~\citep{anderson2004integrated} and Soar~\citep{laird2012soar} consolidate experience through specific mechanisms (chunks, reinforcement) but keep memory stores representationally separate. LLM-based agents~\citep{yao2023react,park2023generative,shinn2023reflexion} typically save experience as text or vectors and reuse it through prompting; consolidation is delegated to the model and is not auditable. A graph offers a middle path: if episodic, semantic, and procedural content share one explicitly typed graph, the write-back of experience---two named layers, continuous Hebbian reinforcement and thresholded structural promotion (terminology note in Section~\ref{sec:arch-wb})---becomes a first-class, auditable operation rather than a side effect of prompting.

This paper reports \textsc{Fascinator} (FAS), a Python prototype built around that idea: a \emph{unified cognitive graph} on which spreading activation~\citep{collins1975spreading} drives attention and candidate generation, an attribution pipeline writes experience back as typed relational structure, and ledger and event-framework channels route stored knowledge into behavior. We deliberately keep two propositions separate throughout:

\begin{itemize}
\item[(a)] \textbf{Persistent relational reuse has behavioral value.} Under information asymmetry---the knowledge needed for a decision exists only in memory, not in the current observation---graph context built from experience changes behavior (beacon task: 20/20 completions vs 0/20 without context).
\item[(b)] \textbf{The necessity of activation is not supported.} The same benefit is fully reproduced by activation-free retrieval (D-flat, 20/20), and a preregistered mechanism campaign rejects activation's incremental value over flat retrieval under all four confirmatory conditions.
\end{itemize}

The contributions are:
\begin{enumerate}
\item \textbf{An auditable unified-graph memory architecture.} Semantic, episodic, and procedural content coexist in one typed weighted graph; experience becomes new relational structure only through a single evidence-weighted write-back channel. Isolation experiments verify the channel (0 vs 13 learned nodes under an identical experience stream), cross-context re-activation (4/12 vs 0/8 vs 0/0 marked nodes), and two consumer channels that measurably flip action selection (Sections~\ref{sec:arch}--\ref{sec:res-persist}).
\item \textbf{A preregistered evaluation protocol with honestly reported negative results.} The campaign testing whether spreading activation adds value over flat retrieval was designed, seeded, and committed before execution; its primary hypothesis was rejected ($\Delta$JCG${<}0$ under all four preregistered conditions, Holm $p{=}1.5\times10^{-9}$, 0/200 reversals), and we report this as a primary result (Section~\ref{sec:res-boundary}).
\item \textbf{A boundary analysis of activation-based access.} The single systematic exception---multi-input convergence through a shared hub (30/30 vs 0/30)---is delimited by an 84-cell parameter sweep over edge weight, input count, and path depth, and complemented by a structural diagnosis of relation-direction semantics (Section~\ref{sec:res-boundary}, \ref{sec:discussion}).
\end{enumerate}

Every number in this paper is programmatically extracted from archived run outputs; the evidentiary standard, the deterministic-sandbox protocol, and the statistics are specified in Section~\ref{sec:method}. Internal codenames, per-family provenance, and full per-experiment data are in the Supplementary Material.

\section{Related Work}
\label{sec:related}

\textbf{Cognitive architectures.} ACT-R~\citep{anderson2004integrated} and Soar~\citep{laird2012soar} treat consolidation as mechanism-specific learning over separate memory stores; OpenCog~\citep{hart2008opencog} shares our atom-space motivation but targets general intelligence rather than auditable consolidation of embodied experience. Episodic-memory theory~\citep{tulving1983elements,baddeley2000episodic} and event segmentation~\citep{zacks2007event,schacter2007remembering} motivate our event-framework channel and attribution window. Spreading activation as an access mechanism goes back to~\citet{collins1975spreading}; attention and working-set selection relate to~\citet{itti2001computational,posner1980orienting,cowan2001magical}.

\textbf{LLM-agent memory.} ReAct~\citep{yao2023react}, Generative Agents~\citep{park2023generative}, Reflexion~\citep{shinn2023reflexion}, and RAG~\citep{lewis2020rag} reuse experience through prompting or vector retrieval; consolidation quality is not separately controllable or auditable. FAS uses the LLM only as a budgeted interface for language and never as memory (Section~\ref{sec:arch-llm}); the studied knowledge channels run with the LLM hard-disabled.

\textbf{Relation to our contribution.} None of these lines reports a preregistered test of whether activation-based access beats flat retrieval on the same structure. Our negative result on that question, and the parameter sweep delimiting the exception, are the paper's methodological core.

\section{Materials and Methods}
\label{sec:arch}

We present the system content-first: the graph, activation, write-back, persistence and consumers, then the LLM boundary.

\subsection{The unified cognitive graph}
\label{sec:arch-graph}
FAS maintains a directed, typed, weighted graph $G=(V,E,W)$. Nodes carry a label $\tau_v$ marking memory role (seven classes---items, recipes, action experiences, changes, events, intentions, gaps---sharing one namespace), a base weight $w_v$, a timestamp, and a runtime activation $a_v\in\mathbb{R}$; each edge $e=(u,v,\rho,w)$ carries a relation type $\rho$ (eight classes: produces, belongs-to, drops, needs, needs-tool, leads-to, involves, results) and a static weight $w\in[-2,2]$; negative weights denote suppression. \emph{There is no separate episodic, skill, or intention store}: goals, events, action patterns, and affect are nodes of the same graph, distinguished only by label and relation type (Fig.~\ref{fig:ontology}). Beside the graph there are only two auxiliary layers---a short-term episodic buffer (capacity 100, for short-horizon integration) and a statistical timeline (the ledger)---which never build structure: the sole structure-writing channel is the attribution pipeline of Section~\ref{sec:arch-wb}. The space split serves lifecycle management only; it does not isolate topology.

\begin{figure}[t]
\centering
\resizebox{\textwidth}{!}{%
\begin{tikzpicture}[
  sp/.style={draw, rounded corners, align=center, inner sep=3.5pt, font=\scriptsize},
  sem/.style={sp, fill=blue!8}, epi/.style={sp, fill=orange!12},
  proc/.style={sp, fill=green!10}, selfst/.style={sp, fill=purple!8}, emo/.style={sp, fill=red!8},
  edgeT/.style={font=\tiny, fill=white, inner sep=1pt},
  typed/.style={-{Latex}, semithick, gray!70}, neg/.style={-{Latex}, thick, red!70, dashed},
  cross/.style={-{Latex}, semithick, blue!60!black},
  corridor/.style={-{Latex}, semithick, blue!60!black, rounded corners=8pt}]
% central graph node
\node[sp, minimum width=4.8cm, minimum height=1.35cm, fill=gray!6, font=\footnotesize, draw=gray!60]
  (G) at (0,0) {Unified cognitive graph $G{=}(V,E,W)$\\ \scriptsize single namespace; shared labels \& relation types\\ \scriptsize \emph{cross-space links are the core capability}};
% semantic cluster (left column)
\node[sem] (rc) at (-4.3,3.4) {recipe\\ \texttt{oak\_planks}};
\node[sem] (it) at (-4.4,2.3) {item\\ \texttt{oak\_log}, \texttt{stick}};
\node[sem] (ch) at (-5.8,1.0) {change\\ \texttt{inventory:oak\_log:}\\ \texttt{count\_increased}};
% episodic (top-right / right)
\node[epi] (ev) at (1.7,3.1) {event\\ \texttt{tree disappeared}};
\node[epi] (ae) at (4.9,1.9) {action experience\\ \texttt{gather(oak\_log)}};
% bottom row
\node[proc] (ac) at (3.0,-2.5) {action skill\\ \texttt{gather\_resource},\\ \texttt{craft\_item}};
\node[selfst] (ci) at (-1.5,-3.0) {intention/gap\\ \texttt{CI\_*}, \texttt{gap:use}};
\node[emo] (em) at (-3.9,-3.2) {affect\\ arousal, stress};
% intra-semantic
\draw[typed] (it) -- node[edgeT, pos=0.30, right]{belongs} (rc);
% cross-space links, corridor-routed (never through G)
\draw[corridor] (rc.north) -- (-4.3,4.15) -- (3.25,4.15) -- (3.25,-1.55) -- (ac.north);
\node[edgeT] at (0.2,4.32){needs};
\draw[corridor] (ae.south) -- (4.9,-4.15) -- (-5.2,-4.15) -- (ch.south);
\node[edgeT] at (-0.4,-4.45){leads-to $w{=}conf$};
\draw[corridor] (it.west) -- (-7.35,2.3) -- (-7.35,-2.3) -- (ac.west);
\node[edgeT] at (-4.0,-2.08){drops};
\draw[cross] (ci) -- node[edgeT, pos=0.3, left]{points-to} (it);
\draw[cross] (ev) -- node[edgeT, pos=0.55, above]{involves} (it);
\draw[cross] (ae) -- node[edgeT, pos=0.72, below]{involves} (it);
\draw[neg] (ev) to[bend left=15] (it);
\node[edgeT, text=red!70] at (-1.3,3.42) {suppress $-w$};
\draw[cross] (em.north) -- node[edgeT, pos=0.45, left]{modulates} (-1.9,-0.72);
\draw[cross] (ac.north) -- node[edgeT, pos=0.5, right]{skill} (2.35,-0.72);
\begin{scope}[on background layer]
\node[draw=blue!30, rounded corners, fit=(it)(rc)(ch), inner sep=3pt] (fsem) {};
\node[draw=orange!40, rounded corners, fit=(ev)(ae), inner sep=3pt] (fepi) {};
\node[draw=green!50, rounded corners, fit=(ac), inner sep=3pt] (fproc) {};
\node[draw=purple!40, rounded corners, fit=(ci), inner sep=3pt] (fself) {};
\end{scope}
\node[font=\scriptsize, blue!60!black, anchor=north east] at ([xshift=-2pt, yshift=-3pt]fsem.south east) {semantic};
\node[font=\scriptsize, orange!70!black, anchor=north] at ([yshift=-3pt]fepi.south) {episodic space};
\node[font=\scriptsize, green!50!black, anchor=north] at ([yshift=-3pt]fproc.south) {procedural space};
\node[font=\scriptsize, purple!60!black, anchor=north] at ([yshift=-3pt]fself.south) {self/intention space};
\end{tikzpicture}}%
\caption{The unified graph ontology (schematic). Seven label classes and eight relation types share one namespace; there is no separate episodic, skill, or intention store. Cross-space links (blue, corridor-routed) are the core capability: the \texttt{leads-to} edge from \texttt{action experience} to \texttt{change} (weight = confidence) is the product of the write-back pipeline (Section~\ref{sec:arch-wb}); \texttt{gap} points to an entity to inject activation; \texttt{event} carries a negative-polarity slot edge (red dashed) suppressing the entity (Section~\ref{sec:arch-persist}).}
\label{fig:ontology}
\end{figure}

\subsection{Cognitive state and spreading activation}
\label{sec:arch-act}
The cognitive state at time $t$ is the activation field $\mathbf{a}(t)$. Inputs (observations, intentions, gap injections) deposit initial energy into the corresponding nodes; each step, every active node $u$ emits to its neighbors
\begin{equation}
\Delta a_{u\rightarrow v}(t) \;=\; a_u(t)\,\cdot\, \tilde w_{uv}\,\cdot\, \beta_{\mathrm{cat}}(\rho)\,\cdot\, g_{\mathrm{rel}}(\rho)\,\Big/\; \textstyle\sum_{v'} \max\big(\tilde w_{uv'},0\big),
\label{eq:spread}
\end{equation}
followed by decay $a_v(t{+}1) = (1-\lambda)\,a_v(t) + \sum_u \Delta a_{u\rightarrow v}(t)$, clamped by per-space maxima (episodic space $a_{\max}{=}3$, global default $5$). Emission follows an ``emit = transfer'' policy: each node's emitted quantity per step (default 1.0) is deducted from itself; source nodes are exempt. This is an engineering choice of short memory, not an energy-conservation law. Negative-weight edges contribute suppressive flow; suppression is issued only when the emitting node's positive-weight total is greater than zero. The emission budget is scaled by a modulatory gain
\begin{equation}
g \;=\; \mathrm{clip}\big(1 + 0.3\,\mathrm{arousal} - 0.15\,\mathrm{stress},\; 0.7,\; 1.6\big),
\label{eq:mod}
\end{equation}
where arousal is derived from affective-node activation and stress from a slow tonic offset. Attention is the Top-$k$ of the field and constitutes input to action selection. Fig.~\ref{fig:act-schematic} sketches activation-based access \emph{side by side with its flat-retrieval comparison}; the schematic builds intuition only---all conclusions are carried by the experiments of Section~\ref{sec:res-boundary}.

\begin{figure}[t]
\centering
\resizebox{\textwidth}{!}{%
\begin{tikzpicture}[
  nd/.style={circle, draw, minimum size=10mm, inner sep=0pt, align=center, font=\scriptsize},
  hub/.style={nd, fill=purple!25}, hot/.style={nd, fill=red!50},
  warm/.style={nd, fill=orange!40}, cool/.style={nd, fill=blue!12},
  annot/.style={align=center, text=black!80, font=\scriptsize},
  flow/.style={-{Latex}, semithick}, flat/.style={-{Latex}, thick, gray!50, dashed},
  inib/.style={-{Latex}, thick, red!70, dashed}]
\node[hot] (a) at (0,2.2) {gap A\\ demand 1}; \node[warm] (x) at (2.0,1.8) {entity a};
\node[hot] (b) at (0,0.0) {gap B\\ demand 2}; \node[warm] (y) at (2.0,0.4) {entity b};
\node[hub] (P) at (3.7,1.1) {shared\\ hub P};
\node[cool] (T) at (5.9,2.0) {target\\ oak\_log};
\node[cool] (rp) at (5.9,0.2) {recipe\\ planks};
\node[cool, draw=gray!60, fill=gray!12] (ev) at (7.7,2.9) {event\\ failure};
\draw[flow] (a) -- (x); \draw[flow] (x) -- (P);
\draw[flow] (b) -- (y); \draw[flow] (y) -- (P);
\draw[flow] (P) -- (rp); \draw[flow] (rp) -- (T);
\draw[inib] (ev) -- (T);
\draw[flat] (-1.35,1.75) -- (-1.35,-0.55) -- (2.5,-0.85) -- (5.35,-0.25);
\node[annot, anchor=east] at (-1.55,1.85) {flat: nearest};
\draw[flat] (0.62,-0.45) -- (2.62,-0.45) -- (4.35,-0.45) -- (5.35,1.60) node[annot, pos=0.20, above, sloped]{Top-1};
\node[annot, align=left, anchor=west] at (0.15,-2.05) {
  Diffusion paths: multi-source input converges via relational structure on hub $P$---downstream convergence (Section~\ref{sec:res-boundary}):\\[-0.5mm]
  \quad $\bullet$ mechanism runs: JCG${>}0$ on 200/200 paired seeds; $\ge$2-input MRR@1${=}1.0$ (hypothesis group)\\[-0.5mm]
  \quad $\bullet$ only systematic positive difference: shared-hub topology, FAS 30/30 \emph{vs} flat 0/30 (single weight $w{=}0.5$ design)\\[-0.5mm]
  Flat path: nearest-neighbor retrieval reaches the target directly---sufficient for ordinary retrieval:\\[-0.5mm]
  \quad $\bullet$ $\Delta$JCG${<}0$ under all 4 preregistered conditions (reverse-direction Holm $p{=}1.5\times10^{-9}$, 0/200 reversals)\\[-0.5mm]
  \quad $\bullet$ wider spread dilutes targets ($\rho{=}{-}0.445$, ${-}0.656$); \emph{not drawn as diffusion-always-wins}.};
\end{tikzpicture}}%
\caption{Spreading activation and flat retrieval side by side. Every numbered claim comes from the archived experiments of Section~\ref{sec:res-boundary}; the figure only builds intuition.}
\label{fig:act-schematic}
\end{figure}

\subsection{Experience attribution and write-back}
\label{sec:arch-wb}
After an action settles, its receipt and the attribution window (8--120\,s of environment time) of observations form candidate action--observation pairs. Repeated pairs aggregate into hypotheses carrying support $s$, contra-evidence $c$, and Laplace confidence $\mathrm{conf}=s/(s+c+2)$. When $s\ge 5$ and $\mathrm{conf}\ge 0.80$, the hypothesis is \emph{promoted} to a typed relation on the graph: a \texttt{leads-to} edge of weight $w=\mathrm{conf}$ from the action node $\mathtt{action:}A(x)$ to the change node $\mathtt{change:}O$,
\begin{equation}
\mathtt{action:}A(x) \;\longrightarrow\; \mathtt{change:}O \qquad (\texttt{leads-to},\; w{=}\mathrm{conf}),
\label{eq:wb}
\end{equation}
and both endpoints are mark-activated. This is how experience becomes \emph{new relational structure}: \texttt{action:}/\texttt{change:} nodes exist only as products of this pipeline. \emph{Terminology note.} This paper distinguishes three operations: \emph{continuous reinforcement} (sub-threshold Hebbian co-activation weight adjustment; mechanism measurements in Sections~\ref{sec:res-psa} and \ref{sec:res-system}), \emph{episodic replay consolidation} (reorganization driven by offline replay, implemented in the intention/replay layer but not behaviorally validated here), and \emph{structural promotion} (the thresholded hypothesis $\rightarrow$ graph write defined in this section). The blanket term ``consolidation'' used in the introduction and related work names the problem-level ideal only; every claim about FAS below is made under the name---and within the evidence scope---of one of these three operations. The same ledger maintains per-action success-rate memories, consumed by action selection as prior bonuses ($\pm0.05$)---the second statistical path from experience to behavior. A third channel encodes discrete events as event-framework nodes with polarity-slot edges (implemented and production-default off; this paper assembles it explicitly, Section~\ref{sec:res-persist}). The algorithm is Algorithm~\ref{alg:wb}.

\begin{algorithm}[t]
\caption{Attribution--write-back: experience to relational knowledge}\label{alg:wb}
\begin{algorithmic}[1]
\Require action--observation stream $(a_t,o_t)$ within the window; ledger $\mathcal{L}$; graph $G$
\State For each action $A_t$, associate observations $O_k$ in the window with its settlement as candidate pairs $(A_t,O_k)$
\State Aggregate by action signature: $s\mathrel{+}=1$ (positive outcome), $c\mathrel{+}=1$ (contra)
\State $\mathrm{conf}=s/(s+c+2)$ (Laplace smoothing)
\If{$s \ge 5$ and $\mathrm{conf} \ge 0.80$}
  \State create \texttt{action:}A, \texttt{change:}O nodes if absent
  \State write edge \texttt{action} $\longrightarrow$ \texttt{change} with relation \texttt{leads-to}, weight $w{=}\mathrm{conf}$; mark both active
\EndIf
\State ledger in parallel maintains per-action success priors $\mathcal{L}.\mathrm{prior}(\mathrm{sig}(A))$ consumed by action selection
\end{algorithmic}
\end{algorithm}

\subsection{Persistence and consumers}
\label{sec:arch-persist}
Knowledge survives episode boundaries by construction (the graph is inherited, not reinitialized); persistence carries two payloads: \emph{structure} and \emph{activation values} are both serialized into the next episode, so ``re-activation across contexts'' is a measurable event. Section~\ref{sec:res-persist} measures that re-activated ($a\ge0.05$) differences are carried exactly by episodic structure connecting shared entities (4/12 vs 0/8 vs 0/0). Snapshot rotation and a shrink guard (save refusal to overwrite a running graph with an empty one) are data-integrity mechanisms, not cognitive content.

Four classes of consumer use the knowledge produced by write-back. \textbf{(a) Ledger channel (production):} candidate action scores are reweighted by ledger success priors ($+0.05$ success / $-0.05$ failure). \textbf{(b) Affect/neuromodulation (production, every tick):} arousal and stress scalars modulate the emission gain via Eq.~\ref{eq:mod}. \textbf{(c) Continuous intentions and replay gating (production):} the continuous-cognition cycle runs graph dynamics every 2.5\,s; intention lifecycle and idle-gated offline experience replay belong here. \textbf{(d) Event framework (experimental assembly, production-default off):} negative events suppress entity activation via polarity-slot edges; the code is production-hardened but disabled by default so that nothing lands on the graph in production; all behavioral evidence comes from experimental assembly. \emph{Experimental layers} (sandbox assembly and the LLM decision head of Section~\ref{sec:method-routing}) never enter the production default path; the run recorder only observes---it never modifies cognitive behavior.

\subsection{The LLM as a budgeted interface}
\label{sec:arch-llm}
FAS uses the LLM for language production, interpretation of unstructured input, and (low-frequency) reflection; it never uses the LLM as memory, nor as the source of the knowledge this paper studies. All subject experiments run under an import-time guard that hard-disables the LLM (zero-LLM contract; any call errors out)---the only exceptions are the explicitly labeled decision-head experiments (Section~\ref{sec:method-routing}).

\subsection{The experience-to-knowledge loop}
\label{sec:arch-loop}
Fig.~\ref{fig:loop} shows the complete conceptual loop. Experience is segmented into temporally ordered action--observation fragments (the attribution window), repeated patterns accumulate evidence into candidate relations, and, after support and confidence checks, are written back as persistent graph structure. New inputs unfold across this graph via spreading activation, forming a task-relevant cognitive state with a Top-$k$ focus; consumers carry that state into action selection; actions change the environment, producing new experience. Language production, interpretation, and reflection are the only LLM calls in the loop (Section~\ref{sec:arch-llm}): the LLM is a consumer/decision component, not memory.

\begin{figure}[t]
\centering
\resizebox{\textwidth}{!}{%
\begin{tikzpicture}[
  loop/.style={draw, rounded corners, align=center, inner sep=4pt, font=\footnotesize},
  scaff/.style={draw, rounded corners, align=center, inner sep=3pt, font=\scriptsize, fill=gray!12, dashed},
  flow/.style={-{Latex}, semithick},
  node distance=3.5mm and 6mm]
\node[loop] (env) {embodied interaction\\ \scriptsize \strut actions, observations\strut};
\node[loop, right=of env] (per) {perception intake\\ \scriptsize \strut ingestion, segmentation\strut};
\node[loop, right=of per] (epi) {action--observation fragments\\ \scriptsize \strut attribution window 8--120 s\strut};
\node[loop, right=of epi] (ev) {evidence accumulation\\ \scriptsize \strut support / contra / confidence\strut};
\node[loop, right=of ev] (wb) {relational write-back\\ \scriptsize \strut promotion to graph + mark-active\strut};
\node[loop, below=of wb] (graph) {unified cognitive graph $G{=}(V,E,W)$\\ \scriptsize \strut structure + activation values persist\strut};
\node[loop, below=of graph] (act) {spreading activation\\ \scriptsize \strut decay, gain, suppression, Top-$k$\strut};
\node[loop, below=of act] (cs) {task-relevant cognitive state\\ \scriptsize \strut working set of relevant knowledge\strut};
\node[loop, below=of cs] (beh) {behavior / action selection\\ \scriptsize \strut consumers: ledger, events, affect\strut};
\node[scaff, below left=6mm and 2mm of ev] (s1) {scaffolding: recipe/block tables};
\node[scaff, left=of cs] (s2) {consumers: reward ledger, trial ledger,\\ event framework (production default off)};
\draw[flow] (env) -- (per);
\draw[flow] (per) -- (epi);
\draw[flow] (epi) -- (ev);
\draw[flow] (ev) -- (wb);
\draw[flow] (wb) -- (graph);
\draw[flow] (graph) -- (act);
\draw[flow] (act) -- (cs);
\draw[flow] (cs) -- (beh);
\draw[flow] (beh) -| ([xshift=-6mm]env.south) -- (env.south);
\draw[flow] (s1) -| (ev);
\draw[flow] (s2) -- ([yshift=-2.5pt]cs.west);
\end{tikzpicture}}%
\caption{The FAS conceptual loop: from embodied interaction to reusable relational knowledge, and back to behavior via activation and consumers. Dashed blocks are scaffolding (mechanics tables; the necessary precondition characterized by the zero-prior result, Section~\ref{sec:res-system}) and consumption channels. The LLM's only positions in the loop are language production/interpretation/reflection (Section~\ref{sec:arch-llm}).}
\label{fig:loop}
\end{figure}

\subsection{Evaluation protocol}
\label{sec:method}

\subsubsection{Environment, determinism, and what ``paired seeds'' mean}
\label{sec:method-env}
All experiments run in a deterministic \texttt{SandboxWorld}: a discrete Minecraft-style grid world with a simulated clock (16\,s/tick), recipe and block tables, and a fake device bridge replacing a real Minecraft LAN bridge. The agent stack is the production stack---knowledge graph, diffusion engine, continuous-cognition components, autonomy loop, action manager, and the real skill layer---with the LLM disabled by the import-time guard. One long-chain condition uses a documented simplification (sandbox smelting settles instantly into the inventory); results relying on it are marked.

Every condition is instantiated with a fixed seed; for the zero-LLM subject experiments, identical seeds reproduce bit-identical trajectories (verified by re-runs). For the LLM decision-head families, see the determinism re-run in Section~\ref{sec:res-rerun}. \emph{Paired seeds} instantiate different random instances of the same condition and support paired statistics across instances; because the world is a deterministic sandbox with synthetic topologies, they do \emph{not} sample real-world variance. Statistical tests in this paper therefore quantify the consistency of a direction of difference across instances, not population-level effect sizes. Sample sizes: subject experiments $n{=}5$ seeds; mechanism campaign $n{=}50$ paired seeds per condition; beacon and routing $n{=}20$ paired seeds; relation-direction diagnosis $n{=}10$ paired seeds. Every run is archived by the run recorder (metadata, event stream, start/end graph snapshots, metrics) into content-addressed directories; the archival census is given in the Supplementary Material (Table~S2).

\subsubsection{Preregistration and evidence accounting}
\label{sec:method-prereg}
The mechanism campaign's hypotheses, topologies, seed reservoir, metrics, and testing scheme were fixed and committed before execution; deviations (one assembly repair before unblinding) are reported as such. Evidence is reported at three levels---\emph{representational} (existence and re-activation of learned nodes/edges), \emph{mechanism} (activation trajectories, weight changes, ledger counts), and \emph{behavioral} (completion ticks, action counts, first-action identity)---and attributed to one of four protocol layers (production default; production mechanism under assembled validation; mechanism-isolation layer on reference models; experimental-layer harness). Table~\ref{tab:evidence} merges the two scales. General-capability claims (``L4'') are explicitly not made. Where the word ``causal'' is used, it means rule-like regularities inferred from temporally ordered action--outcome observations with evidence weights---not interventional identification; no do-style test is performed anywhere. \emph{Statistical conventions.} Every comparison family declares its multiplicity correction in its own section: the beacon uses exact McNemar with no multiplicity; the mechanism campaign uses a one-sided normal-approximation Wilcoxon (rank-based, analytic tail) with Holm correction over its four confirmatory conditions (direction preregistered as FAS${>}$B2; the frozen implementation defines the one-sided $p$ as half the two-sided $p$ regardless of sign, so with the all-negative $\Delta$JCG observed here the reported Holm value quantifies the reverse-direction tail, and the preregistered-direction tail is its complement, ${\approx}1$); the relation-direction diagnosis is a single declared family of 21 tests; the controls experiment applies Bonferroni $m{=}8$; the persistent-state experiment tests three preregistered hypotheses; campaign v2 applies Holm within each experiment. One-sided tests are used only where the preregistration declares them (C-1--C-4; persistent-state H5; E2/T); all other reported tests are two-sided, and exact methods are used at $n{\le}25$. Table~
ef{tab:stats} lists every family with its test, sidedness, pairing, and correction (full version with implementation provenance: Supplementary Material~S18).

\begin{table}[t]
\centering\small
\caption{Evidence levels $\times$ protocol layers (merged scale; used to attribute every claim in this paper).}
\label{tab:evidence}
\begin{tabular}{L{3.2cm}L{4.4cm}L{6.4cm}}
\toprule
 & \textbf{Level of the claim} & \textbf{Layer the evidence comes from} \\
\midrule
L1 representational & learned structure exists and re-activates & production default path (zero-LLM subject behavior) \\
L2 mechanism & measurable activation/ledger dynamics & production mechanism under experimental assembly; mechanism-isolation layer (reference models distilled from production semantics) \\
L3 behavioral & stored knowledge changes action selection & assembled consumer channels; decision-head harness (experimental layer, never attached to production) \\
L4 general capability & --- & \emph{explicitly not claimed anywhere in this paper} \\
\bottomrule
\end{tabular}
\end{table}

\subsubsection{The mechanism-campaign protocol (preregistered)}
\label{sec:method-civ}
The core incremental-value campaign studies the \emph{incremental} value of spreading activation over flat retrieval. Four synthetic subject topologies are used (3 inputs/100 nodes/25\% density; 3 inputs/300 nodes/50\%; 5 inputs/300 nodes/50\%; 5 inputs/1000 nodes/75\%). Four candidate generators are compared: \textbf{B1} direct injection (graph-free); \textbf{B2} \emph{flat path-product}---a closed-form approximation of retrieval on the same structure (per input--candidate pair, the best simple-path edge-weight product, summed over inputs; same structural information, no iterative propagation; note this is a closed-form surrogate, not an embedding retriever---the beacon task's flat control in Section~\ref{sec:res-reuse} is a genuine sentence-embedding retriever); \textbf{B3} random activation (equal-energy control); \textbf{B4} oracle (ground truth, used only for generator-correctness prechecks; MRR${=}1.0$ on all four topologies). The primary metric JCG (joint convergence gain) measures multi-input convergence onto shared targets. The preregistered confirmatory conditions C-1--C-4 ($n{=}50$ paired seeds each) use one-sided paired Wilcoxon with Holm correction to test H3 (\emph{diffusion's JCG exceeds flat's}). Ancillary families (all run and archived): insertion-order permutations, label/ID permutation (hidden oracle), a no-joint-target counterfactual, a negative-edge test, a weight sweep $w\in\{0.1..0.9\}$, noise/distractor sweeps, and combinational-topology detectors. Condition-level numbers are in the Supplementary Material (Table~S6); the preregistration ships with the artifacts.

\subsubsection{LLM decision-head protocol (routing experiments)}
\label{sec:method-routing}
One family places a real LLM at the decision point (single model, temperature 0, one call per decision). Menu, environment state, task facts, experience stream, and per-decision budgets are identical across conditions; the only difference is the context the model sees: (A) none; (B) raw last-8 experiences; (C) embedding-similarity retrieval top-5; (D) FAS mechanism-selected context (graph entry, propagation, Top-$k$ focus, plus cognitive demand/gap/routing annotations); (E) random memory entries (counterfactual control). Four task classes are used, all fully observable at decision time. The repaired official campaign runs on the production edge-building path with nine fail-fast harness invariants (520/520 runs valid). The campaign is positioned as exploratory: it measures context-selection value, and the FAS loop does not control execution.

\subsubsection{Reproducibility}
Fixed seeds; single-process deterministic execution; no wall-clock-dependent logic in the measurement path (one documented exception above); archived raw runs and content summaries; all regression checks pass at the time of writing: 82 unit tests in the pytest suite and 108/108 checks in the regression-script runner. A determinism re-run check for the LLM decision-head family is reported in Section~\ref{sec:res-rerun}.

\begin{table}[t]
\centering\footnotesize
\caption{Statistical conventions, one row per comparison family (full version with implementation locations in Supplementary Material~S18). ``Status'': pre $=$ preregistered; post $=$ post-hoc; desc $=$ descriptive (no test). All $p$ values in this paper are produced by the test named in this table.}
\label{tab:stats}
\begin{tabular}{L{3.0cm}L{3.4cm}L{2.2cm}L{1.1cm}L{2.2cm}L{0.9cm}}
\toprule
\textbf{Family} & \textbf{Test} & \textbf{Sidedness} & \textbf{paired $n$} & \textbf{Correction ($m$)} & \textbf{Status} \\
\midrule
Beacon success (\S5.1) & exact McNemar & two-sided & 20 & none & pre \\
Beacon decisions (\S5.1) & paired Wilcoxon & two-sided & 20 & none & pre \\
Mechanism $\Delta$JCG, C-1--C-4 (\S5.3) & Wilcoxon, normal approx.\ (analytic tail) & one-sided; reported Holm value is the \emph{reverse}-direction tail (preregistered-direction tail ${\approx}1$) & 50 & Holm (4) & pre \\
Mechanism secondary (MRR, JCG${>}0$) (\S5.3) & paired Wilcoxon & two-sided & 50 & none (secondary family) & pre \\
Shared-hub sweep (\S5.3) & exact McNemar; Wilcoxon & two-sided & 30 & none (frozen design) & pre \\
Relation-direction diagnosis (\S5.3) & Wilcoxon (exact) & two-sided & 10 & single declared family (21 tests) & pre \\
Campaign v2 (C/E2/T/I/B) (\S5.5) & McNemar; Wilcoxon; descriptive counts & two-sided; one-sided for E2/T; B descriptive & 10 (T: 9) & Holm (within experiment) & pre \\
Controls C3/C4 (\S5.6) & exact McNemar; Wilcoxon & two-sided & 10 & Bonferroni (8) & pre \\
Persistent state H1/H3 (\S5.6) & exact McNemar & two-sided & 10 & none (three hypotheses) & pre \\
Persistent state H5 (\S5.6) & Wilcoxon & one-sided & 10 & none & pre \\
Routing (\S5.7) & paired Wilcoxon & two-sided & 20 & Bonferroni (280) & post \\
Cross-episode recombination (\S5.7) & McNemar; Wilcoxon & two-sided; one-sided (greater) & 30 & within experiment & pre \\
Determinism re-run (\S5.4) & agreement counts & --- & 20 & --- & desc \\
Presentation-equivalence revision (\S5.8) & none (stopped at gate; 0 runs) & --- & 0 & --- & post \\
\bottomrule
\end{tabular}
\end{table}

\section{Results}
\label{sec:results}

\subsection{Persistent relational reuse: the information-asymmetric beacon task}
\label{sec:res-reuse}

\emph{Task.} The beacon task operationalizes information asymmetry: the knowledge needed to solve the task exists \emph{only} in experience-written memory, in no input prompt. The target recipe is stripped from all task facts (``closed book''); the crafting menu lists \texttt{craft\_X} only while item X's context is present (menu gating); goal entities are excluded from activation seeding. Four conditions, $n{=}20$ paired seeds, temperature 0, all 80 runs valid: \textbf{fas\_full} (persistent graph + activation context), \textbf{D-noact} (same graph context with activation dynamics disabled), \textbf{D-flat} (flat retrieval without activation aggregation), \textbf{llm\_direct} (no memory context).

\begin{table}[t]
\centering\small
\caption{Beacon task results ($n{=}20$ paired seeds per condition; 80 runs). All numbers from the archived campaign analysis.}
\label{tab:beacon}
\begin{tabular}{L{2.6cm}L{2.2cm}L{3.4cm}L{3.4cm}L{2.6cm}}
\toprule
\textbf{Condition} & \textbf{Completed} & \textbf{Decisions (mean, median)} & \textbf{First action \texttt{gather\_oak\_log}} & \textbf{Contrast} \\
\midrule
fas\_full & 20/20 & 2.15, 2 & 20/20 & vs direct: McNemar $p{=}2\times10^{-6}$ \\
D-flat (embedding retrieval) & 20/20 & 2.1, 2 & 19/20 & vs fas: $p{=}1.0$; decisions Wilcoxon $p{=}0.72$ \\
D-noact & 0/20 & 8.0, 8 & 20/20 & vs fas: McNemar $p{=}2\times10^{-6}$ \\
llm\_direct & 0/20 & 8.0, 8 & 17/20 & --- \\
\bottomrule
\end{tabular}
\end{table}

\emph{Results.} fas\_full completes 20/20; llm\_direct completes 0/20 (McNemar exact $p{=}2\times10^{-6}$; every decision window is used up without unlocking the chain). \textbf{D-noact}---structure present, runtime access absent---fails 0/20: structure alone is not usable access. \textbf{D-flat}---activation-free flat aggregation---fully reproduces the benefit (20/20; decisions mean 2.1 vs 2.15, no statistical difference, Wilcoxon $p{=}0.72$). First-action identity gives a mechanism fingerprint: fas and D-noact both choose \texttt{gather\_oak\_log} 20/20, direct only 17/20---the difference is not in first-action choice but in subsequent unlocking.

Two readings must be kept separate. \textbf{(a) Persistent relational reuse produces behavioral benefit} on a high-noise, recipe-opaque, information-asymmetric task. The 0/20 baselines are, by task design, an \emph{existence verification}: closed-book gating makes the recipe unavailable to context-free conditions by construction, so the fas-vs-direct contrast demonstrates that the pipeline delivers the required knowledge through memory, not that memory context is ``statistically significantly better'' in any population sense. \textbf{(b) The necessity of activation is not supported}: D-flat's complete reproduction joins the mechanism campaign of Section~\ref{sec:res-boundary}---flat retrieval suffices on reachable tasks.

\emph{What kind of retriever is D-flat?} D-flat is a genuine sentence-embedding retriever: it encodes all graph node identifiers with a public, locally runnable sentence-transformer model (BAAI/bge-small-zh-v1.5) and ranks them by cosine similarity to the task text---no activation, no propagation, no path structure. A zero-LLM accessibility check (Supplementary Material~S6) verifies the enabling condition directly: for the beacon goal text, the recipe node \texttt{item:oak\_planks} ranks 2nd and the resource node \texttt{item:oak\_log} 4th among all 44 graph nodes, inside the Top-8 context window, for all 20 seed assemblies. A complete real case study of the wood-to-plank chain (practice experience, write-back, persistent graph, beacon behavior) is shown in Supplementary Figure~S3.

\subsection{Write-back isolation, persistence, and consumption}
\label{sec:res-persist}

\emph{Write-back isolation.} If written-back structure were a rearrangement of injected knowledge, disabling write-back would change nothing. The isolation experiment breaks graph write-back on a single switch: perception, actions, and environment are identical; two lineages practice the same four-step chain for eight rounds in the same world. With write-back off, the final graph contains \textbf{0} learned \texttt{action:}/\texttt{change:} nodes; with write-back on, \textbf{13} (Supplementary Figure~S1a). Sorting the 15 promoted relations of the write-back-on lineage (Supplementary Table~S4): 11 substantive action--effect associations (``mining yields items and removes a block''), 4 degenerate relations (action pointing at its own success event---the promotion rules do not distinguish this class; we report rather than hide it), and 2 substantive pairs polluted by the attribution window (downstream craft effects attributed to the upstream gather). The initial recipe tables know \emph{what can craft what}; \emph{what execution produces} can only be established by doing.

\emph{Cross-context re-activation.} Episode 1 practices the chain in one context; Episode 2 changes the surface context completely (interfering stand, different layout) with a structurally similar goal. Three conditions inherit from Episode 1: full graph (G0), flat-memory control with episodic-space nodes stripped (G1), and a clean graph (G2). Across 5/5 seeds: G0 re-activates \textbf{4 of 12} marked nodes ($a\ge0.05$); G1 \textbf{0 of 8}; G2 has none (Supplementary Figure~S2). The re-activation difference is carried exactly by the episodic event structure connecting shared entities; behavioral metrics are identical across conditions---representational/access-layer differences are separated from the behavior layer.

\emph{Ledger consumption.} After 8 practice rounds of a three-step chain (ledger promotes 7 relations), the agent migrates to a structurally related four-step task with the ledger inherited (Transfer) or fresh (No-Transfer). Deterministic, 5/5 seeds: Transfer completes in 13 ticks / 8 actions; No-Transfer needs 16 ticks / 9 actions; the first action flips from a curiosity-driven detour to the experienced resource, and the difference lands exactly on the three candidate actions shared by both tasks (prior $(1.0,\ \text{obs}{=}8)$ \emph{vs} (None, 0)). Task success is 5/5 in both conditions (Supplementary Figure~S4).

\emph{Event-framework consumption.} A negative event (a tree that disappears on first dig) is encoded as an event-framework node whose negative-polarity slot edge suppresses entity activation. After three warm-up ticks, the suppressed curiosity path loses the first-action competition: 5/5 control first-actions at tick 9 \emph{vs} inhibited action at tick 6 (3 ticks earlier, 1 curious action saved). With \emph{real} failure receipts (written only when an action genuinely fails), control retries the doomed target once (n$=$4, first completion attempt @17), the event-framework condition never (n$=$3, @15)---and because that failure mode is a transient timeout, the statistical channel is fully silent on it: the graph-side event framework is the only trace of the failure and the only source of the behavioral difference. Scope: this channel is assembled at the experimental layer and production default is off. It is evidence that graph-encoded experience \emph{can} flip action selection, not evidence that production does so today.

\subsection{Preregistered boundary analysis of activation-based access}
\label{sec:res-boundary}

\emph{The convergence mechanism genuinely runs.} The mechanism campaign first asks whether multi-input convergence on the graph is a real phenomenon: JCG${>}$0 on \textbf{200/200} paired seeds (each confirmatory condition $n{=}50$; median JCG 0.0977/0.0977/0.2158/0.2177); the hypothesis group ($\ge$2 inputs) reaches MRR@1${=}1.0$. Counterfactual controls behave as required: with no joint target, all four methods score 0.0 (convergence is not a self-fulfilling prophecy); under label/ID permutation (hidden oracle, 30 seeds), target rank 22.5 is unchanged (no structural privilege); insertion order (10 permutations) gives CV${=}0.0$; the negative-edge test leaves totals unchanged (0.0333${=}$0.0333; total 17.41). An independent 663-run falsification campaign had first confirmed the convergence reference model (JCG${>}0$ rate 0.970). At the access layer, structural diffusion is also distinguishable from noise and from lookup: on a four-condition access task, full diffusion produces activation quality 61.5 over a field of 39 lit nodes, diffusion-off 16.6 over 17, the equal-energy random control 45.8 over a larger structureless field, and a target-chain oracle 17.2 over 18 (Supplementary Figure~S5).

\emph{Incremental value over flat retrieval: rejected.} Against the closed-form flat path-product (B2), all four preregistered conditions show $\Delta$JCG${<}0$ (median differences ${-}0.402$, ${-}0.402$, ${-}0.784$, ${-}0.782$; paired $n{=}50$ per condition; the frozen implementation's one-sided $p$ is half the two-sided $p$ regardless of direction, so on these all-negative data the reported tail is the \emph{reverse} direction (B2${>}$FAS): $p{=}3.8\times10^{-10}$, Holm ($m{=}4$) $1.5\times10^{-9}$, while the preregistered-direction (FAS${>}$B2) tail is its complement, $p{\approx}1$; effect $r{=}0.615$; seeds where FAS beats flat: \textbf{0/200}). MRR: FAS below B2 (0.0435/0.0426 vs 0.5/0.5 on the two smaller topologies; tie at 1.0 on the two larger). Noise robustness is unsupported: robustness-sweep slope 0.0076 (B2 0.0), AUC 0.0427 vs 0.04. \textbf{H3 is rejected}: under these topologies and information budgets, flat retrieval is no worse than diffusion, and usually better (Table~\ref{tab:boundary}, Fig.~\ref{fig:boundary}a).

\begin{table}[t]
\centering\small
\caption{Preregistered boundary analysis: key mechanism-level numbers. All values are read programmatically from the archived statistics files; full data in Supplementary Material S7--S9.}
\label{tab:boundary}
\begin{tabular}{L{4.4cm}L{6.6cm}L{3.4cm}}
\toprule
\textbf{Test} & \textbf{Numbers} & \textbf{Reading} \\
\midrule
JCG${>}0$ (mechanism runs) & 200/200 paired seeds; medians 0.0977/0.0977/0.2158/0.2177 & convergence mechanism is real \\
$\Delta$JCG (C-1--C-4, preregistered) & ${-}0.402$/${-}0.402$/${-}0.784$/${-}0.782$; reverse-direction Holm $p{=}1.5\times10^{-9}$, $r{=}0.615$, 0/200 reversals & \textbf{H3 rejected} \\
MRR (FAS vs B2, smaller topologies) & 0.0435/0.0426 vs 0.5/0.5 & flat better on reachable tasks \\
Noise robustness & slope 0.0076 vs 0.0; AUC 0.0427 vs 0.04 & no diffusion advantage \\
Shared-hub topology (single weight $w{=}0.5$) & FAS 30/30 vs B2 0/30, B1 0/30, B3 5/30 & the one systematic exception \\
Weight sweep & both methods reach MRR${=}1.0$ together at $w{=}0.7$ & structural threshold, not a FAS advantage \\
Dilution & $\rho(\mathrm{reach},\mathrm{act})$: $-0.445$ (noise), $-0.656$ (distractor) & wider spread, flatter target \\
Shared-hub parameter sweep (84 cells $\times$ 30 seeds) & FAS strictly better in 11 cells (all $w{\le}0.5$, $d{\le}1$); flat strictly better in 9 (all $w{\ge}0.7$, $d{\ge}1$); tie in 64; O3 replication 30/30 vs 0/30 & advantage region: weak edges $\times$ shallow convergence paths \\
\bottomrule
\end{tabular}
\end{table}

\emph{The shared-hub exception, parameterized.} The only systematic positive difference in the preregistered campaign appears in the combinational topology: multiple inputs' relevant paths converge through a \emph{shared hub} (FAS 30/30; B2 0/30, B1 0/30, B3 5/30; single weight configuration $w{=}0.5$). Because this exception carried the paper's only positive evidence for activation, we ran a preregistered-style parameter sweep (design frozen before execution; Supplementary Material~S8): joint grid over hub-path edge weight $w\in\{0.3,\ldots,0.9\}$, number of converging inputs $k\in\{2,3,4,6\}$, and private-path depth $d\in\{0,1,2\}$, 30 paired seeds per cell, FAS vs the same closed-form flat baseline, with the original single-hub configuration as a hard preflight (replicated: FAS 30/30 vs B2 0/30). The result is a mapped surface, not a point (Fig.~\ref{fig:sweep}): FAS is strictly better in 11 of 84 cells---\emph{all} with weak weights ($w{\le}0.5$) and shallow convergence paths ($d{\le}1$); flat is strictly better in 9 cells---\emph{all} with strong weights ($w{\ge}0.7$) and deeper paths ($d{\ge}1$); the remaining 64 tie. On the O3-shaped axis ($k{=}2$, $d{=}0$), flat catches up exactly when the hub-path weight reaches 0.6 (both methods 30/30 at $w{\ge}0.6$). The mechanism reading is consistent with the dilution result: convergence helps when the hub can accumulate multi-input energy faster than single paths saturate their targets, and hurts when strong edges let flat ranking resolve the target directly or when long weak paths dilute the spread.

\begin{figure}[t]
\centering
\includegraphics[width=\textwidth]{figures/fig_sweep.pdf}
\caption{Shared-hub parameter sweep (84 cells $\times$ 30 paired seeds, deterministic, zero-LLM). Cell color: paired success difference (FAS minus flat) on the joint target; green = FAS strictly better, orange = flat strictly better, gray = tie; annotations give the cell counts. The FAS-advantage region is confined to weak edge weights and shallow convergence paths. All numbers read programmatically from the archived sweep outputs.}
\label{fig:sweep}
\end{figure}

\emph{Weight threshold, not diffusion advantage.} In the campaign's weight sweep, both methods jump to MRR${=}1.0$ together at $w{=}0.7$ (FAS 0.083@0.5${\rightarrow}$1.0@0.7; B2 0.077${\rightarrow}$1.0)---a structural resolvability threshold, not a FAS-specific advantage (Fig.~\ref{fig:boundary}b). The dilution trade-off completes the picture: the wider the diffusion front (more reachable nodes), the lower the target activation ($\rho{=}{-}0.445$ noise sweep, ${-}0.656$ distractor sweep; Fig.~\ref{fig:boundary}c).

\begin{figure}[t]
\centering
\includegraphics[width=\textwidth]{figures/fig_diffusion_boundary.pdf}
\caption{The diffusion applicability boundary. (a) Preregistered confirmatory family: all four $\Delta$JCG negative (reverse-direction Holm $p{=}1.5\times10^{-9}$); the only positive point is the shared-hub topology (star, single-weight $w{=}0.5$ design). (b) Weight sweep: both methods reach MRR${=}1.0$ together at $w{=}0.7$. (c) Dilution trade-off: the wider the spread, the flatter the target activation ($\rho{=}{-}0.445$, ${-}0.656$).}
\label{fig:boundary}
\end{figure}

\emph{Relation-direction semantics (diagnostic).} A separate mechanism diagnosis (10 paired seeds; full data in Supplementary Material~S9) showed that the production recipe graph is a depth-1 star forest whose relation targets are all sinks, making forward two-hop propagation structurally impossible; adding generic reverse-edge projections restores multi-hop reachability and roughly doubles the rank correlation between topological score and activation ($\rho$ 0.325${\rightarrow}$0.645), at the cost of diluting single-node activation (0.833${\rightarrow}$0.475). All 21 paired tests give $p{=}0.001953$---the minimum attainable two-sided $p$ at $n{=}10$ (exact Wilcoxon, all seeds concordant), which indicates uniform direction and nothing more. We report this as a design-level diagnostic (the fix is an ontology-design decision, and the production graph was not modified), not as behavioral evidence.

\subsection{Determinism re-run of the decision-head family}
\label{sec:res-rerun}
Because the beacon and routing families use a real LLM decision head (temperature 0), we re-ran a 5-seed $\times$ 4-condition subset of the beacon campaign in a fresh output directory and compared against the archived runs. The result is an honest negative: \textbf{task outcomes} (success, per condition) match in 20/20 re-runs, but \textbf{full action sequences} match bit-for-bit in only 7/20---the remote endpoint is not bit-deterministic even at temperature 0. We therefore claim trajectory-level determinism only for the zero-LLM families (bit-identical by construction and verified by re-runs); the decision-head families are outcome-robust but not trajectory-reproducible, and no claim in this paper rests on their exact trajectories (Supplementary Material~S10).

\subsection{Comparative behavioral evidence (campaign v2)}
\label{sec:res-campaign}

The protocol above tests activation against closed-form and embedding
retrieval on \emph{mechanism} metrics. To test the same questions on
\emph{behavior}, we ran a second campaign (508 archived runs; the 390
runs from the two mid-course confound repairs are archived
separately; design and full data in Supplementary
Material~S15--S17) in which every condition shares the same scripted
experience stream, the same menu, the same LLM decision head
(temperature 0), and the same decision budget---only the memory form
differs.

\textbf{Goal resumption under neutral instructions (C).} After an
open-book practice phase (obtain \texttt{oak\_planks}) and an unrelated
distractor phase, the test prompt gives no goal at all
(``you have free choice'') and the recipe facts are stripped. FAS
resumes and completes the abandoned goal in \textbf{9/10} paired seeds;
the full-history LLM baseline and the no-memory baseline resume in
\textbf{0/10} (McNemar $p{=}0.0039$; Holm $p{=}0.0117$). This advantage
is regime-dependent: with a longer experience stream that also contains
a failure episode, the full-history baseline matches FAS at 10/10
(below)---the salience advantage holds in the short-history regime.
Mechanism attribution: a zero-experience control reproduces the
same working set (the recipe-closure prior edge
\texttt{oak\_log}$\rightarrow$\texttt{oak\_planks} re-lights the chain
when the raw material is observed), so the effect is
\emph{observation-triggered activation over prior relational
structure}---relational working-set salience---rather than an
intention-lifecycle mechanism.

\textbf{Goal revision is where the same mechanism costs (E).} With an
infeasible target (tool-gated ore), the full-history baseline stops
repeating the failed gather after \textbf{1.0} failures on average and
pivots immediately; FAS repeats it \textbf{10.5} times and pivots in
7/10 runs (late). A joint measurement (below) gives the direction and
magnitude of this cost precisely. The two results are the same coin: state that keeps
a goal salient under distraction also keeps it salient when the
environment invalidates it.

\textbf{Consumer repairs (E2/F).} We then tested whether the failure
and goal states are simply not consumed by the decision context.
Wiring the production event-frame's negative-polarity encoding into
the campaign assembly (failure$\rightarrow$entity edge $w{=}{-}0.8$,
parameters identical to the production config) reduced repeated
failures from 10.5 to \textbf{7.2} (one-sided exact Wilcoxon
$p{=}0.0098$)---statistically significant, but still far weaker than the history baseline's 1.0 repeats, confirming residual design-level goal pressure,
so residual inertia is design-level (constant goal text, weak ledger
prior). Adding a uniform self-graph goal consumer
(\texttt{current\_goal} $\rightarrow$ working set; production default
off) turned the previously bit-identical ON/OFF behavior into a
separation: with the registry cleared, the consumer-on loop re-derives
the abandoned goal in 5/5 seeds and its gathering actions rise from
1.0 to \textbf{6.0} per run. The self-model's ``behavioral
irrelevance'' was a missing-consumer artifact, not a property of the
representation.

\textbf{Behavioral convergence (I).} On a two-chain convergence task
(\texttt{oak\_fence} requires \texttt{oak\_planks} and \texttt{stick}
learned in separate experiences), the main batch (per condition
$n{=}10$) shows FAS experience-subset conditions at 0/10 (I1--I3) and
2/10 when both streams are supplied together (I4)---hence an overall
0--2/10---while the raw-history LLM baseline completes \textbf{10/10}
and the retrieval baseline 8/10 (I3 \emph{vs} I5, Holm $p{=}0.041$).
In the representation batch (30 runs), context budget ($k{=}16,32$)
and a structure-enriched serialization (explicit
action$\rightarrow$outcome pairs) both remain at 0/10, locating the
failure in the \emph{serialization format}: the graph state drops the
temporal narrative that the raw stream preserves (Supplementary
Material~S17).

\textbf{Credit horizon (B) and long-run stability (J).} A deterministic
delay probe on the production attribution pipeline (mechanism batch:
delays $0/32/80/160$\,s, 10 seeds each, 40 runs) shows inventory-side
action$\rightarrow$consequence credit forming in 10/10 at every delay
up to 80\,s and in 0/10 at 160\,s (the self-success channel is already
0/10 at 80\,s)---the 90\,s attribution window is the
architecture's credit horizon (an intentional semantic boundary,
Section~\ref{sec:limits}). A 100-cycle stability run (zero-LLM) shows
bounded graph growth ($+3.7\%$ nodes, $245{\rightarrow}254$), clamped
activation, and no contamination collapse, but also zero promotion
(repetition-starved episodes never reach the support threshold);
short-range within-task Hebbian reinforcement of the ten prior-closure
edges has been naturally reproduced in the persistent-state experiment
(Section~\ref{sec:res-psa}), but \emph{long-run} (hundred-cycle,
cross-episode) continuous reinforcement remains untested.

\textbf{Joint measurement (T).} On a single task sequence carrying
both mechanisms (a failed infeasible goal followed by a neutral free
phase), the event-frame encoding replicates its effect
(failed-gather repetitions 12.0$\rightarrow$7.8, one-sided exact
Wilcoxon $p{=}0.002$, $n{=}9$ valid pairs), the goal consumer adds no resumption increment where the
prior-edge path is already at ceiling (all conditions 10/10) and adds
no inertia penalty (6.8 vs 7.8, n.s.)---the two consumers coexist, and
the combined working set crowds the failed action out of the top-8
entirely. A corrected resumption probe (fresh world) shows all
conditions---including the full-history baseline---at 10/10, confirming
that the resumption advantage is regime-dependent and that this
comparison does not isolate graph structure or activation.
The behavioral trade-off curve between suppression and
salience is therefore measurable, and its first two points favor
suppression. (One design flaw---inventory carry-over saturating the
first resumption metric---was caught and corrected with a fresh-world
rerun; both datasets are archived.)

Together these results revise the paper's claim structure: persistent
relational state demonstrably changes behavior in one direction
(resumption), demonstrably costs in another (revision), is currently
bottlenecked at the serialization interface (convergence), and is
credit-limited by an explicit design window (delayed consequence).
All negative and reverse-direction results are reported at full
strength.

\subsection{Persistent-state test under context removal}
\label{sec:res-psa}

The campaign above varies only the memory form; the persistent-state
experiment removes the \emph{experience stream itself} from the
decision context to test the behavioral role of persistence directly
(preregistered; design, all conditions, and data in Supplementary
Material~S16). Three tasks: T1 hunger recovery (existence check), T2
goal resumption under context removal (practice a goal, then a fresh
context with recipe facts stripped and only graph-derived content in
the decision context), and T3 failure repetition after the environment
invalidates the goal. Conditions: \textbf{FAS-full} (persistent graph
context), \textbf{FAS-reset} (same stack, graph and activation reset at
the phase boundary), \textbf{B0-direct} (no context), budgeted raw
history \textbf{h128/h512/h2048}, \textbf{B2-rag} (embedding
retrieval), and \textbf{FAS-sham} (node names without experience
edges). Each hypothesis test uses $n{=}10$ paired seeds; same LLM
decision head (temperature 0), same menu and decision budget.

\textbf{Results.} On T2, FAS-full achieves \textbf{10/10}, FAS-reset
\textbf{0/10} and B0-direct \textbf{0/10} (H1/H3: exact McNemar
$p{=}0.00195$---the minimum attainable at $n{=}10$ with all pairs
discordant, evidencing directional consistency only, not effect size).
H5 (T3): mean failure repetitions \textbf{2.8} (FAS-full) \emph{vs}
\textbf{7.7} (FAS-reset), one-sided Wilcoxon $p{=}0.00098$;
B0-direct 11.2; h2048 and B2-rag both 0.0. Mechanism observation: the
ten prior-closure edges underwent Hebbian continuous reinforcement
during phase~1---a short-range, within-task natural reproduction
(contrasting with the 30-episode whitelisted drift of
Section~\ref{sec:res-system}).

\textbf{What this experiment does not show.} (i) FAS-reset and
B0-direct are \emph{empty controls} (state cleared / context absent):
the experiment establishes that persistence \emph{as a whole} is
causally operative, not the independent contribution of graph
structure or activation. The structure-preserving no-activation
controls (flat store, activation-off) complete T2 at 10/10
each---\emph{no difference detected} against FAS-full ($n{=}10$,
limited power)---and their worse T3 behavior (12.0/11.7 repeats) is
traceable to failure-evidence visibility: the control implementations
structurally exclude failure information from the decision context.
A frozen-preregistration revision attempting presentation equivalence
was gate-stopped before any run (Section~\ref{sec:audit-negative}); no
presentation-equivalent graph control exists yet. (ii) FAS-sham also reaches
\textbf{10/10}: node-name traces plus menu gating already suffice at
this task surface. (iii) The budget curve is non-diagnostic: the
experience stream is only 465 tokens and observation lines embed
inventory snapshots, leaking answer-bearing strings into truncated
history---h128/h512/h2048 score 10/10, 10/10, 9/10 and B2-rag 10/10:
whenever the experience stream stays re-readable, history and
retrieval baselines match FAS (H2 does not hold). (iv) Single task
family, $n{=}10$; all significance statements are directional claims.

\subsection{Additional observations}
\label{sec:res-system}
Five further families are compressed here; full data in Supplementary Material~S11--S12.

\emph{Long chain and priors.} A 14-step chain completes 5/5 seeds at 315 ticks; removing graph prior edges \emph{and} the smelting task prior gives 0/5 (the agent cannot plan out-of-recipe transitions); re-adding only the smelting prior restores the bit-identical 315 ticks; $-$diffusion is faster (259 ticks). With zero priors, the system never enters a learning regime: 0/25 episodes, the goal system locks onto an unreachable target, and the attribution ledger is completely untouched (0 aggregations) while the graph still grows from perception ($+29$ nodes)---\emph{scaffolded grounded knowledge acquisition}: experience does not automatically become knowledge.

\emph{Hebbian, affect, intentions.} Hebbian co-activation strengthening changes exactly 3 whitelisted edges ($+0.074$ total on the dominant edge; total-sum drift $+0.079$) over 30 episodes with retrieval ranks bitwise unchanged---mechanism isolation without behavioral effect. Scalar affective modulation moves activation quality monotonically with the gain (49.8 / 52.5 / 57.4 for gains 0.85/1.0/1.3; disabled is bitwise identical to neutral) but moves no action-level metric. The continuous-intention lifecycle forms, competes, and decays from real experience (3--4 intentions per run); the re-fire gate never triggers (field temperature 1.61--2.12 vs cold-field threshold 0.35), and the expression level is never reached (peak 0.46 $<$ threshold 0.78)---formation through decay is validated, expression is not.

\emph{Explicit routing (exploratory).} On four fully-observable decision-head task classes (520/520 valid runs), FAS context shows one directional positive (competing-targets task: 19/20 vs 15/20 without context, uncorrected $p{=}0.046$; after correction over the campaign's 280 tests, $p{\approx}1$), one efficiency deficit (re-routing task: 4.15 vs 2.3/2.8 decisions, $p{=}0.003$/$0.044$), and structural token costs (2.1--3.7$\times$). An earlier routing campaign was rejected by our own assembly audit (the harness never built edges and swallowed 521/521 exceptions) and is archived as invalid. Under this paper's standard, routing sits with the negative results.

\emph{Live observations.} Two 25-minute live Minecraft sessions (zero LLM in the action path) are observational only: the first grows the graph by $+100$ nodes/$+365$ edges in discovery form; the second adds only $+35$/$+95$ with behavior dominated by learned goals; the live causal ledger (89 aggregations / 78 hypotheses / 29 promotions) comes entirely from real behavior. The causal bridge from write-back to behavioral improvement is \emph{not} established in the live window---an honest negative observation.

\emph{Closure masking.} Across many families, completion-layer metrics are systematically insensitive to internal differences (write-back lineages, re-activation conditions, and modulation conditions all complete identically). On tasks whose closure guarantees reachability, the planner alone suffices and activation-level knowledge falls outside the decision envelope; evaluations that report only closure-guaranteed completion rates will systematically fail to detect memory phenomena.

\subsection{Audit-driven corrections and negative results}
\label{sec:audit-negative}

Four methodological corrections emerged during this work and are
reported as first-class contributions:

\textbf{(C) Misattributed resumption.} The goal-resumption advantage
(9/10 vs 0/10, Section~\ref{sec:res-campaign}) was initially labelled
``persistent intention.'' A zero-experience control reproduces the same
working set: the driver is the prior-closure edge re-lit by observation,
not an intention-lifecycle mechanism. The effect is relational
working-set salience, regime-dependent (short-history only).

\textbf{(E) Assembly defect.} The goal-revision deficit (10.5 failed
repeats vs history's 1.0) was traced to a missing consumer: the
production event-frame's negative-polarity channel was not wired into
the campaign assembly. Wiring it reduced repeats to 7.2
(one-sided exact $p{=}0.0098$)---a partial repair, with residual inertia attributable
to design-level goal pressure.

\textbf{(F) Consumer gap.} The self-model's goal state was written to
the graph (telemetry 5/5/5) but had zero production readers. Adding a
uniform state consumer changed autonomous behavior (gathers 1.0
$\rightarrow$ 6.0), confirming the write-only hypothesis.

\textbf{(T3) Presentation inequivalence.} The T3 advantage of FAS-full
over flat-store and no-activation controls was traced to
failure-evidence visibility: the control implementations structurally
excluded failure information from the decision context. A
frozen-preregistration revision (flat failure-first, 2-hop adjacency,
2-hop without polarity edges) was attempted and stopped at its own
presentation-equivalence gate, before any LLM run: the flat
failure-first control passed the gate, but under the naive 2-hop rule
the failure-evidence node is structurally reachable (within a 51-node
2-hop set) yet ranked 43rd by discovery-edge weight (0.5 vs a slot-8
weight of 0.8) and truncated out of the 8-slot budget. Failure-evidence
absence in graph controls is therefore jointly determined by distance,
ranking weight, and budget---not hop distance alone---and information
equivalence does not imply presentation equivalence. The preregistered
stop rule took effect as designed; the independent contribution of
graph structure and activation remains untested.

\textbf{Negative results preserved:} behavioral convergence (I) fails
at all budgets; the credit horizon (B) is 90\,s; cross-episode
behavioral advantage over retrieval does not hold; history/RAG
baselines match FAS whenever the experience stream is short enough to
fit their budget.

\section{Discussion}
\label{sec:discussion}

\subsection{What FAS demonstrates}
\textbf{Shown, within the exact scope of this protocol:} (i) under mechanism-level scaffolding, embodied experience converts into relational structure absent from the initial graph and all injected sources (0 vs 13 nodes; 15 promotions); (ii) that structure persists and re-activates in changed contexts (4/12 vs 0/8 vs 0/0); (iii) two named consumers (ledger, event framework) measurably change behavior (13 vs 16 ticks; 5/5 first-action flips); (iv) under information asymmetry, persistent relational reuse yields a large behavioral difference (20/20 vs 0/20), fully reproducible by activation-free embedding retrieval; (v) the convergence mechanism runs (200/200) yet its incremental value over flat retrieval is preregistration-rejected (0/200 reversals), with the shared-hub exception delimited by the 84-cell sweep; (vi) under context removal, persistence as a whole is causally operative on behavior (FAS-full \emph{vs} the graph-reset empty control: T2 10/10 \emph{vs} 0/10, exact McNemar $p{=}0.00195$; T3 failure repeats 2.8 \emph{vs} 7.7, one-sided Wilcoxon $p{=}0.00098$)---shown only against empty controls. \textbf{Not shown:} general intelligence, open-ended learning from scratch, universal superiority over LLM memory systems, causal understanding, universal diffusion benefits, robust routing value, a behavioral Hebbian benefit, or the independent contribution of graph structure and activation over structure-preserving controls (untested: no presentation-equivalent graph control has been realized; Section~\ref{sec:audit-negative}). Each item is accompanied by its evidence or its absence.

\subsection{Why flat retrieval usually suffices---and when it does not}
The negative results organize into a testable conditional statement rather than a slogan. When the task only asks ``retrieve the related thing,'' flat retrieval does the same work at lower complexity (beacon: D-flat 20/20; campaign: $\Delta$JCG${<}0$). The mechanism-level reason is dilution: the wider the diffusion front, the lower the target activation ($\rho{=}{-}0.445$, ${-}0.656$), which also explains why $-$diffusion is faster on the long chain (259 vs 315 ticks). The sweep sharpens the exception into a region: \emph{activation-based access pays when the task requires multi-relation convergence and the relevant edges are too weak for direct ranking to resolve the target} (all 11 FAS-advantage cells have $w{\le}0.5$, $d{\le}1$); with strong edges, flat ranking resolves everything; with deep weak paths, diffusion dilutes faster than the closed-form product decays (all 9 flat-advantage cells). The guarded hypothesis is: \emph{when decision-relevant information is absent from current observation and the task structure requires multi-relation convergence, activation-based access has positive value over flat retrieval; when information is observable or the task is nearest-neighbor, that value is zero or negative.} Four evidence lines meet at this hypothesis---the beacon (information asymmetry), the shared-hub cells (convergence), the rejected campaign (nearest-neighbor), and the routing tasks (observable information)---none of them stamps it alone.

\subsection{The role of scaffolding}
The zero-prior result (0/25 with an untouched ledger) shows scaffolding is not a convenience but a precondition: the goal system locks onto an unreachable target before any attributable experience exists, and the learning channel starves. FAS is therefore \emph{scaffolded grounded} rather than \emph{learning from scratch}: the architecture must provide the mechanisms (attribution window, thresholds, mechanics tables) that make experience readable and writable. This holds for any system that must act before it can learn, and it is both a necessary transformation and a boundary of capability.

\subsection{Why the unified graph matters}
The three functional evidence chains (re-activation carried by episodic structure, exact ledger transfer onto shared candidates, suppression flips through polarity-slot edges) exist because experience, semantics, and action candidates \emph{share one addressable structure}: in separated stores, each coupling would require custom bridging; here they are natural consequences of graph traversal and graph-based activation. The value of the unified graph is of the form \emph{observable consequences}, and the negative results (closure masking, sub-threshold Hebbian drift) delimit what it cannot do: representational changes do not automatically become behavioral changes.

\subsection{Design implication: direction semantics}
The relation-direction diagnostic (Section~\ref{sec:res-boundary}; full data in Supplementary Material~S9) generalizes into an ontology-design rule: \emph{bake per-relation-class direction semantics into the ontology rather than assume propagation will find its way}---recipe-class relations forward, observation-co-occurrence relations bidirectional. On the measured recipe graph (44 nodes/53 edges), forward two-hop reachability is structurally zero because every relation target is a sink; the improvement after adding reverse projections is a design-level verification, not an emergent property.

\subsection{Limitations}
\label{sec:limits}
(1) All experiments run in a deterministic sandbox with synthetic topologies; paired seeds instantiate different random instances, not real-world variance, and all significance statements quantify directional consistency, not population effects. (2) The beacon's 0/20 baselines are existence verifications built into the task design (closed-book gating), not independent statistical findings. (3) Single environment family, single Minecraft-style domain; no third-party benchmarks; no external-environment validation (e.g., a Crafter-style equivalent) was performed. (4) Subject experiments use five deterministic seeds---descriptive statistics only; paired experiments ($n{=}10$, $n{=}20$, $n{=}50$) carry the directional claims with stated corrections. (5) Semantic write-back structure is not consumed by the planner; learned knowledge affects behavior through the ledger, event-framework, and re-activation channels only. (6) Zero-prior conditions test locking, not learning ability. (7) Not implemented: prospective memory, full perception pipeline, symbolic reasoning, RL policy nodes; several implemented mechanisms are not experimentally validated here. (8) Event-framework behavioral results come from experimental assembly of a production-default-off mechanism. (9) Scaffolding dependence is not quantified. (10) Experimental graphs are small (44--1000 nodes); scale limits all mechanism conclusions. (11) Routing tasks are fully observable at decision time and the negative results are conditional on that; token costs and directional differences are exploratory. (12) The mechanism campaign and relation-direction diagnosis are mechanism-level evidence without production behavioral tests; the shared-hub sweep, while preregistered-style and frozen before execution, is still confined to synthetic hub topologies with a closed-form flat surrogate, not an embedding retriever. (13) The beacon is a single task family (craft domain, closed book); cross-task generality of the information-asymmetry benefit is untested. (14) The decision-head families depend on a remote LLM endpoint at temperature 0; bit-level determinism is verified only by re-running (Section~\ref{sec:res-rerun}) (15) The attribution window (90\,s) is a hard credit horizon: delayed consequences beyond it are never assigned (Section~\ref{sec:res-campaign}). (16) The decision-head context serialization drops the temporal narrative of the experience stream; behavioral convergence tasks therefore fail at this interface. (17) The resumption advantage over full-history baselines is regime-dependent (short-history regime only). (18) The persistent-state experiment (Section~\ref{sec:res-psa}; FAS-full vs FAS-reset, 10/0) used a single task family with a 465-token experience stream; its primary McNemar $p{=}0.00195$ is the minimum attainable at $n{=}10$ and reflects directional consistency, not effect size; the budget curve is non-diagnostic because inventory snapshots in the observation lines leak answer-bearing strings into truncated history. (19) FAS-sham (nodes without experience edges) achieves 10/10 on the goal task, indicating that node-name traces plus menu gating suffice at this task surface. (20) On goal completion (T2), no difference is detected between FAS and the flat-store/no-activation controls (10/10 each, $n{=}10$, limited power); FAS-full's T3 advantage over these controls is attributable to failure-evidence visibility (the control implementations exclude failure information from the decision context), so the independent contribution of graph structure or activation is \emph{not} established. A frozen-preregistration revision attempt stopped at its own presentation-equivalence gate (the naive 2-hop rule still truncates the failure evidence), leaving the question untested; the gate also showed that (a) it checked only step-0 contexts over the Phase-1 baseline store, not Phase-3-evolved text, and (b) C1's context carries additional sections (edges/demand/gap/routing; 22 vs 10 lines) that the graph controls lack, so context structure was never fully aligned; answering the question requires a new, independently preregistered design.

\section{Conclusions}
On a unified cognitive graph, with mechanism-level scaffolding, embodied experience becomes new, typed, evidence-weighted relational knowledge: it persists, re-activates across contexts, and reaches behavior through named consumer channels in measurable, deterministic ways. On the information-asymmetric beacon task, persistent relational reuse yields a 20/20 vs 0/20 behavioral difference, fully reproduced without activation. Spreading activation is a real, measurable access mechanism whose incremental value over flat retrieval is preregistration-rejected everywhere except a now-parameterized shared-hub region: weak edges, shallow convergence paths. The context-removal experiment further locates persistence itself as a causal source of behavior (10/0 against the empty control), while the attribution of that effect to graph structure or activation over structure-preserving controls remains open---the audit marks this boundary explicitly, and the gate-stopped revision attempt leaves the question to future, independently preregistered work. The boundaries are equally clear: no-scaffolding learning regimes are never entered; closure-guaranteed tasks mask internal differences; diffusion's footprint is knowledge access, not action selection. We offer the architecture and the protocol---deterministic, zero-LLM, fully archived, with two reported assembly audits and a preregistered rejection---as a substrate for studying when experience becomes knowledge that changes behavior, and we hand over all negative results as part of the map.

\paragraph{Data and code availability.} All experiments are deterministic and fully archived (raw run outputs, graph snapshots, statistics files); a one-click reproduction script, the preregistration records, and the claim-to-evidence mapping ship with the artifacts (Supplementary Material~S14).


% ---------------------------------------------------------------------------
% PeerJ declaration block (all fields are author-confirmation TODOs)
% ---------------------------------------------------------------------------
\section*{Funding}
[TODO-FOR-AUTHOR: state funding or ``The authors received no external funding for this work.'']

\section*{Competing interests}
[TODO-FOR-AUTHOR: state competing interests or ``The authors declare no competing interests.'']

\section*{Author contributions}
[TODO-FOR-AUTHOR: CRediT roles, e.g. ``KY conceived, designed, implemented, and wrote the work.'']

\section*{Data and code availability}
All experiments are deterministic and fully archived; the complete reproducibility package --- raw run outputs, graph snapshots, statistics files, source code, preregistration records, and a one-command verifier that recomputes the paper's key numbers without an LLM --- is openly available on Zenodo at [ZENODO-DOI] under [LICENSE-TO-BE-CHOSEN]. Preregistrations are included in the package (also Supplementary Material S14--S16). LLM decision-head replications require an anonymized endpoint configuration that is not redistributed; they reproduce at the outcome level rather than bit-identical trajectories (Section~\ref{sec:res-rerun}).

\section*{Reporting guidelines}
The study follows preregistered-experiment reporting: preregistration records for the mechanism campaign, the persistent-state experiment and controls, and the presentation-equivalence revision (including its gate-stopped preflight) ship with the reproducibility package and are summarized in Supplementary Material S14--S16. No animal or human subjects were involved; no ARRIVE/PRISMA/STROBE checklist applies. [TODO-FOR-AUTHOR: confirm.]

\bibliographystyle{plainnat}
\bibliography{references}

\end{document}
